\chapter{Algèbres de Lie}\label{ch_liealg}
\minitoc
\newpage

\noindent
Dans tout ce chapitre, soit $K$ un corps.

\section{Généralités}
\subsection{Définitions et premières propriétés}

\defi{
    Une \emphindex{algèbre de Lie} $\mathfrak{g}$ est un $K$-espace vectoriel munie d'une application bilinéaire $\mathfrak{g}\times \mathfrak{g} \to \mathfrak{g}$, notée $(x,y) \mapsto [x, y]$ et appelée \emphindex{crochet de Lie}, vérifiant:
    \ilist{
        \item anticomutativité\footnote{si le corps est de caractéristique différente de $2$, cet axiome est équivalent à demander que $\forall x \in \mathfrak{g}, \quad [x,x] = 0$}: $\forall x, y \in \mathfrak{g}, \quad [x, y] = - [y, x]$
        \item identité de Jacobi\index{identité de Jacobi}: $\forall x, y, z \in \mathfrak{g}, \quad [x, [y, z]] + [y, [z, x]] + [z, [x, y]] = 0$
    }
}

\defi{
    Un \emphindex{morphisme d'algèbres de Lie} est une application linéaire $\varphi: \mathfrak{g} \to \mathfrak{h}$ telle que 
    $$\forall x, y \in \mathfrak{g}, \quad \varphi([x,y]) = [\varphi(x), \varphi(y)]$$
    Un \emph{endomorphisme} est un morphisme $\mathfrak{g} \to \mathfrak{g}$. \\
    Un \emph{isomorphisme} est un morphisme bijectif. \\
    Un \emph{automorphisme} est un morphisme $\mathfrak{g} \to \mathfrak{g}$ bijectif.
}

\defi{
    Une \emphindex{sous-algèbre de Lie} de $\mathfrak{g}$ est un sous-espace vectoriel $\mathfrak{h} \subset \mathfrak{g}$ stable par le crochet de Lie, c'est à dire :
    $$\forall x, y \in \mathfrak{h}, \quad [x,y] \in \mathfrak{h}$$ $\mathfrak{h}$ est donc une algèbre de Lie pour le crochet de Lie induit.
}

\defi{
    Un \emphindex{idéal} de $\mathfrak{g}$ est un sous-espace vectoriel $I \subset \mathfrak{g}$ absorbant pour le crochet de Lie, c'est à dire :
    $$\forall x \in \mathfrak{g}, \forall y \in I, \quad [x, y] \in I$$ 
}

\prop{
    Soit $\varphi: \mathfrak{g} \to \mathfrak{h}$ un morphisme d'algèbres de Lie.
    \ilist{
        \item L'image directe $\Img(\varphi)$ est une sous-algèbre de Lie de $\mathfrak{h}$.
        \item Le noyau $\ker(\varphi)$ est un idéal de $\mathfrak{g}$.
    }
}

%TODO: preuve

\subsection{Algèbre de Lie quotient et théorèmes d'isomorphismes}

\defi{
    Soit $I$ un idéal de $\mathfrak{g}$. L'\emphindex{algèbre de Lie quotient} $\mathfrak{g}/I$ est le $K$-espace vectoriel quotient muni du crochet de Lie
    $$[\bar{x}, \bar{y}] = \overline{[x,y]}$$
    où $\bar{x} = x + I$ désigne la classe de $x$ dans $\mathfrak{g}/I$. Ce crochet est bien défini car $I$ est un idéal.
}


\thm{
    Soit $\varphi : \mathfrak{g} \to \mathfrak{h}$ un morphisme d'algèbres de Lie et $I \subset \ker(\varphi)$ un idéal de $\mathfrak{g}$. Alors il existe un unique morphisme d'algèbres de Lie $\bar{\varphi} : \mathfrak{g}/I \to \mathfrak{h}$ tel que $\varphi = \bar{\varphi} \circ \pi$, où $\pi : \mathfrak{g} \to \mathfrak{g}/I$ est la surjection canonique.
}

%TODO: preuve

\cor{
    Soit $\varphi : \mathfrak{g} \to \mathfrak{h}$ un morphisme d'algèbres de Lie. Alors
    $$\mathfrak{g}/\ker(\varphi) \cong \Img(\varphi)$$
}

%TODO: preuve

\thm{
    Soit $\mathfrak{h}$ une sous-algèbre de Lie de $\mathfrak{g}$ et $I$ un idéal de $\mathfrak{g}$. Alors $\mathfrak{h} + I$ est une sous-algèbre de Lie de $\mathfrak{g}$, $\mathfrak{h} \cap I$ est un idéal de $\mathfrak{h}$, et
    $$\mathfrak{h}/(\mathfrak{h} \cap I) \cong (\mathfrak{h} + I)/I$$
}

%TODO: preuve

\thm{
    Soit $I \subset J$ deux idéaux de $\mathfrak{g}$. Alors $J/I$ est un idéal de $\mathfrak{g}/I$ et
    $$(\mathfrak{g}/I)/(J/I) \cong \mathfrak{g}/J$$
}

%TODO: preuve

\subsection{Somme directe d'algèbres de Lie}

\defi{
    La \emphindex{somme directe} de deux algèbres de Lie $\mathfrak{g}$ et $\mathfrak{h}$ est le $K$-espace vectoriel $\mathfrak{g} \times \mathfrak{h}$ muni du crochet de Lie
    $$[(x_1, y_1), (x_2, y_2)] = ([x_1, x_2]_{\mathfrak{g}},\ [y_1, y_2]_{\mathfrak{h}})$$
    On la note $\mathfrak{g} \oplus \mathfrak{h}$.
}

\prop{
    La somme directe $\mathfrak{g} \oplus \mathfrak{h}$ vérifie les propriétés universelles suivantes : pour toute algèbre de Lie $\mathfrak{k}$,
    \ilist{
        \item pour tout couple de morphismes $\varphi_1 : \mathfrak{k} \to \mathfrak{g}$ et $\varphi_2 : \mathfrak{k} \to \mathfrak{h}$, il existe un unique morphisme $\varphi : \mathfrak{k} \to \mathfrak{g} \oplus \mathfrak{h}$ tel que $\pi_1 \circ \varphi = \varphi_1$ et $\pi_2 \circ \varphi = \varphi_2$
        \item pour tout couple de morphismes $\psi_1 : \mathfrak{g} \to \mathfrak{k}$ et $\psi_2 : \mathfrak{h} \to \mathfrak{k}$, il existe un unique morphisme $\psi : \mathfrak{g} \oplus \mathfrak{h} \to \mathfrak{k}$ tel que $\psi \circ \iota_1 = \psi_1$ et $\psi \circ \iota_2 = \psi_2$
    }
    où $\pi_1, \pi_2$ sont les projections canoniques et $\iota_1, \iota_2$ les injections canoniques.
}

%TODO: preuve

\cor{
    Les sous-espaces $\mathfrak{g} \times \{0\}$ et $\{0\} \times \mathfrak{h}$ sont des idéaux de $\mathfrak{g} \oplus \mathfrak{h}$, isomorphes respectivement à $\mathfrak{g}$ et $\mathfrak{h}$. Réciproquement, si une algèbre de Lie $\mathfrak{k}$ admet deux idéaux $I_1, I_2$ tels que $I_1 \cap I_2 = 0$ et $I_1 + I_2 = \mathfrak{k}$, alors $\mathfrak{k} \cong I_1 \oplus I_2$.
}

%TODO: preuve

\subsection{Algèbres de Lie linéaires}

\defi{
    L'\emphindex{algèbre de Lie générale} est l'algèbre de Lie $\mathfrak{gl}(n, K)$ des matrices carrées $\mathcal{M}(n, K)$ munie du crochet de Lie
    $$[X, Y] = XY - YX$$
}

\defi{
    L'\emphindex{algèbre de Lie spéciale linéaire} est la sous-algèbre de Lie
    $$\mathfrak{sl}(n, K) = \ker(\tr) = \{ X \in \mathfrak{gl}(n, K) \mid \tr(X) = 0 \}$$
}

\defi{
    L'\emphindex{algèbre de Lie symplectique} est la sous-algèbre de Lie
    $$\mathfrak{sp}(2n, K) = \left\{ X \in \mathfrak{gl}(2n, K) \mid \transpose{X} J + J X = 0 \right\}$$
    où $J = \begin{pmatrix} 0 & I_n \\ -I_n & 0 \end{pmatrix}$ est la matrice symplectique standard.
}

\defi{
    L'\emphindex{algèbre de Lie orthogonale} est la sous-algèbre de Lie réelle
    $$\mathfrak{o}(n, \mathbb{R}) = \mathfrak{so}(n, \mathbb{R}) = \{ X \in \mathfrak{gl}(n, \mathbb{R}) \mid \transpose{X} + X = 0 \}$$
    Plus généralement, pour une forme de signature\footnote{La convention choisie ici place les valeurs négatives en premier dans $I_{p,q}$. L'autre convention, qui place les valeurs positives en premier, est également répandue et donne une algèbre de Lie isomorphe.} $(p, q)$, on définit
    $$\mathfrak{o}(p, q, \mathbb{R}) = \mathfrak{so}(p, q, \mathbb{R}) = \{ X \in \mathfrak{gl}(p+q, \mathbb{R}) \mid \transpose{X} I_{p,q} + I_{p,q} X = 0 \}$$
    où $I_{p,q} = \begin{pmatrix} -I_p & 0 \\ 0 & I_q \end{pmatrix}$.
}

\defi{
    L'\emphindex{algèbre de Lie unitaire} est la sous-algèbre de Lie complexe
    $$\mathfrak{u}(n) = \{ X \in \mathfrak{gl}(n, \mathbb{C}) \mid X^\dagger + X = 0 \}$$
    où $X^\dagger = \transpose{{\overline{X}}}$ désigne le conjugué hermitien. L'\emphindex{algèbre de Lie unitaire spéciale} est
    $$\mathfrak{su}(n) = \{ X \in \mathfrak{u}(n) \mid \tr(X) = 0 \}$$
}

\section{Algèbre de Lie nilpotentes}

\subsection{Généralités}

\defi{
    La \emphindex{série centrale descendante} d'une algèbre de Lie $\mathfrak{g}$ est la suite d'idéaux de $\mathfrak{g}$ définie par :
    $$\mathfrak{g}_0 = \mathfrak{g} \quad \text{et} \quad \mathfrak{g}_{i+1} = [\mathfrak{g}, \mathfrak{g}_i]$$
}

\defi{
    Une \emphindex{algèbre de Lie nilpotente} est une algèbre de Lie $\mathfrak{g}$ telle que sa série centrale descendante est nulle à partir d'un certain rang. C'est à dire qu'il existe $i \in \N$ tel que $\mathfrak{g}_i = 0$.
}

%TODO: sous-algèbre, image directe, extensions d'algèbres de Lie nilpotentes sont nilpotentes

\subsection{Théorème d'Engel}

\thm{
    Une algèbre de Lie $\mathfrak{g}$ est nilpotente si et seulement si 
    $$\forall x\in \mathfrak{g}, \quad \ad x \quad \text{est nilpotent}$$ 
}
%TODO: preuve du thèorème d'Engel

\section{Algèbres de Lie résolubles}

\subsection{Généralités}

\defi{
    La \emphindex{série dérivée} d'une algèbre de Lie $\mathfrak{g}$ est la suite d'idéaux de $\mathfrak{g}$ définie par :
    $$\mathfrak{g}^0 = \mathfrak{g} \quad \text{et} \quad \mathfrak{g}^{i+1} = [\mathfrak{g}^i, \mathfrak{g}^i]$$
}

\defi{
    Une \emphindex{algèbre de Lie résoluble} est une algèbre de Lie $\mathfrak{g}$ telle que sa série dérivée est nulle à partir d'un certain rang. C'est à dire qu'il existe $i \in \N$ tel que $\mathfrak{g}^i = 0$.
}

%TODO: sous-algèbre, image directe, extensions d'algèbres résolubles et somme d'idéaux résolubles sont résoluble

%TODO: nilpotent => résoluble

%TODO: définiton du radical d'une algèbre de Lie

%TODO: définition du centre d'une algèbre de Lie

\subsection{Théorème de Lie}

\noindent Dans cette partie on suppose que $K$ est algébriquement clos et de charactéristique nulle.

\thm{
    Soit $\mathfrak{g}$ une sous-algèbre de Lie résoluble de $\mathfrak{gl}(V)$ où $V$ est un $K$-espace vectoriel non nul de dimension finie. Alors $V$ possède un vecteur propre commun à tous les élements de $\mathfrak{g}$
}
%TODO: preuve du théorème de Lie

\subsection{Critère de Cartan}

\thm{
    Soit $\mathfrak{g}$ une sous-algèbre de $\mathfrak{gl}(V)$, où $V$ est un $K$-espace vectoriel de dimension finie, telle que $\tr(xy)=0$ pour tout $x\in [\mathfrak{g}, \mathfrak{g}]$ et tout $y\in\mathfrak{g}$. Alors $\mathfrak{g}$ est résoluble.
}
%TODO: preuve du critère de Cartan

\section{Algèbres de Lie semi-simples}

\defi{
    Une \emphindex{algèbre de Lie simple} est une algèbre de Lie $\mathfrak{g}$ qui ne possède aucun idéal autres que $0$ et elle-même.
}

\defi{
    une \emphindex{algèbre de Lie semi-simple} est une algèbre de Lie $\mathfrak{g}$ telle que $\rad(\mathfrak{g}) = 0$.
}

\subsection{Forme de Killing}

\defi{
    Soit $\mathfrak{g}$ une algèbre de Lie. La \emphindex{forme de Killing} de $\mathfrak{g}$ est l'application $\kappa$ définie par:
    $$\forall x, y\in \mathfrak{g}, \quad \kappa(x,y) = \tr(\ad x \ad y)$$
}

\prop{
    Soit $\mathfrak{g}$ une algèbre de Lie. Alors $\mathfrak{g}$ est semi-simple si et seulement si sa forme de Killing est non-dégénérée.
}

%TODO: semi-simple = somme directe d'idéaux simples

%TODO: semi-simple => toutes les dérivations sont interieures

%TODO: decomposition de Jordan abstraite
